\documentclass[a4paper,fleqn]{cas-dc}

\usepackage[numbers,sort&compress]{natbib}
\usepackage{graphicx}

\begin{document}

\title[mode=title]{Multimodal Fusion for Industrial Anomaly Detection}

\author[1]{Wei Chen}
\cormark[1]
\ead{wei.chen@nthu.edu.tw}

\author[2]{John Smith}

\affiliation[1]{organization={National Tsing Hua University},
  addressline={101 Sec. 2 Kuang-Fu Rd.},
  city={Hsinchu},
  postcode={30013},
  country={Taiwan}}

\affiliation[2]{organization={Massachusetts Institute of Technology},
  addressline={77 Massachusetts Ave.},
  city={Cambridge},
  state={MA},
  postcode={02139},
  country={USA}}

\cortext[1]{Corresponding author}

\begin{abstract}
Industrial anomaly detection has become critical for smart manufacturing,
where early identification of process deviations can prevent costly failures.
We propose a multimodal fusion architecture that combines visual inspection
images with time-series sensor data to detect anomalies in real time.
Building on the attention-based framework of \citep{smith2020}, our method
introduces a cross-modal alignment module that learns shared representations
across heterogeneous sensor streams. The fusion is performed at an intermediate
feature level, preserving modality-specific information while enabling joint
reasoning. We evaluate on three industrial benchmark datasets: MVTec AD,
a printed circuit board inspection dataset we collected in-house, and
a public turbine vibration dataset. Our approach achieves state-of-the-art
image-level AUROC of 98.2 percent on MVTec AD, outperforming the strongest
baseline by 2.1 points. On the PCB dataset the gain is 4.7 points, and on
the turbine dataset we observe consistent improvements across all fault
types. Ablations confirm that both the cross-modal alignment module and
the intermediate-level fusion are essential; removing either drops
performance substantially. We further analyze failure cases and discuss
deployment considerations including model compression and edge inference.
\end{abstract}

\begin{keywords}
anomaly detection \sep multimodal learning \sep smart manufacturing \sep sensor fusion
\end{keywords}

\begin{highlights}
\item A cross-modal alignment module fuses vision and sensor streams
\item Evaluated on three industrial benchmarks including a new PCB dataset
\item State-of-the-art performance on MVTec AD, PCB inspection, and turbine vibration detection benchmarks
\item Ablations confirm both fusion module and alignment are essential
\end{highlights}

\section{Introduction}

Industrial anomaly detection is foundational to smart manufacturing
\citep{smith2020,doe2021}. Our architecture is shown in Figure~\ref{fig:arch}.
The color images are preprocessed before feature extraction.

\section{Related Work}

Prior approaches to multimodal fusion \citep{jones2022,brown2023} have
largely focused on late fusion strategies.

\section{Method}

Figure~\ref{fig:arch} depicts the overall pipeline. The cross-modal
alignment loss is defined as
\begin{equation}
L_\text{align} = \| \phi_v(x_v) - \phi_s(x_s) \|_2^2.
\end{equation}

\begin{figure}
\includegraphics[width=\linewidth]{arch.pdf}
\caption{Architecture overview of the proposed multimodal fusion pipeline.}
\label{fig:arch}
\end{figure}

\section{Experiments}

Results on the three benchmarks are shown in Figure~\ref{fig:results}.

\begin{figure}
\includegraphics[width=\linewidth]{results.pdf}
\caption{Quantitative results on three benchmark datasets.}
\label{fig:results}
\end{figure}

\begin{figure}
\includegraphics[width=\linewidth]{ablation.pdf}
\caption{Ablation study.}
\label{fig:ablation}
\end{figure}

\section{Conclusion}

We analyse the behaviour of our model and find it optimises the cross-modal
objective effectively. We believe this opens promising directions.

\section*{Acknowledgements}

We thank the anonymous reviewers for their helpful comments.

\section*{CRediT authorship contribution statement}

\textbf{Wei Chen:} Conceptualization, Methodology, Software, Writing -- original draft.
\textbf{John Smith:} Supervision, Writing -- review \& editing.

\section*{Declaration of competing interest}

The authors declare that they have no known competing financial interests
or personal relationships that could have appeared to influence the work
reported in this paper.

\bibliographystyle{cas-model2-names}
\bibliography{refs}

\end{document}
